\section{引言}

近年来，深度生成模型的学习重新受到关注 \citep{greedy-dbn,dbm,DARN,vae,vae2}。
多数方法面临的共同困难是训练时需要进行后验推断：大多数潜变量模型的对数似然梯度都由后验统计量定义（例如 \citet{dbm,deep-sigmoid,DARN}）。一种应对方式是在生成模型之外同时训练识别网络 \citep{helmholtz_machine}。给定观测值，识别网络预测潜变量的后验分布，而且通常能以远快于 MCMC 等通用推断算法的速度给出粗略近似。

变分自编码器（VAE；\citet{vae,vae2}）是近期提出的生成模型，它将自顶向下的生成网络与自底向上的识别网络相结合。
两个网络联合训练，以最大化数据对数似然的变分下界。
近期，VAE 已成功用于分离风格与内容 \citep{vae-semi-supervised,tejas_vae}，以及学习以逼真的方式“绘制”图像 \citep{draw}。

VAE 对后验分布作出了较强假设。通常，VAE 假定后验近似按维度因子化，而且其参数可通过对观测值的非线性回归来预测。由于训练目标是最大化对数似然的变分下界，模型会倾向于学习满足这些假设的表示，即后验近似因子化、且能由神经网络预测的表示。这种作用虽有好处，却也有代价：限制后验的形式会限制模型的表达能力。VAE 目标尤其如此：即使识别网络将大部分概率质量放在能够很好解释数据的区域，它仍会严厉惩罚那些难以解释数据的近似后验样本。

本文提出重要性加权自编码器（IWAE）。它与 VAE 采用相同架构，但使用由重要性加权导出的更紧的对数似然下界进行训练。识别网络生成多个近似后验样本，再对其权重取平均。随着样本数增加，在适当的可积性条件下，该下界趋近真实对数似然。使用多个样本使 IWAE 能更灵活地学习后验分布不满足 VAE 建模假设的生成模型。这一方法与重加权唤醒—睡眠算法 \citep{RWS} 有关，但 IWAE 使用单一、统一的目标进行训练。
与 VAE 相比，IWAE 能学习具有更多潜在维度的丰富表示，从而在密度估计基准上获得显著更高的对数似然。
